\documentclass[12pt]{book}
\usepackage{amsmath}
\usepackage{amssymb}
\usepackage{geometry}
\geometry{margin=1in}

\title{Volume 11: Validation, Benchmarks, and Empirical Tests \\ Book 01: Benchmark Architecture}
\author{James C. Harvey}
\date{December 2025}

\begin{document}
\maketitle
\tableofcontents

\chapter{Validation as Geometric Consistency}

\section*{Abstract}
SDT validation is the confirmation of geometric derivations against observed values. This book defines
the benchmark structure and documents the validation philosophy: no fitted constants, only derived
geometry.

\section{Benchmark Pathway}
Each benchmark follows the sequence: definition $\rightarrow$ derivation $\rightarrow$ numerical
evaluation $\rightarrow$ comparison.

\section{Error Structure}
Error is measured relative to derived constants, not adjusted by parameter fitting.

\chapter{Benchmark Taxonomy}

\section*{Abstract}
Benchmarks are grouped by regime: atomic, nuclear, molecular, and cosmological. Each group has distinct
geometry inputs but identical derivation chains.

\section{Regime Groups}
\begin{equation}
\mathcal{B} = \{\mathcal{B}_{\text{atomic}}, \mathcal{B}_{\text{nuclear}}, \mathcal{B}_{\text{molecular}}, \mathcal{B}_{\text{cosmic}}\}.
\end{equation}

\chapter{Reproducibility Protocol}

\section*{Abstract}
Reproducibility is enforced by code and data artifacts. All benchmarks are re-runnable with fixed
constants and documented inputs.

\section{Protocol Steps}
Each benchmark provides: constant set, geometry inputs, derivation sequence, numeric output, and
comparison table.
\end{document}
